\section{Other Arities, Variances, and the Signature Without Top}
\label{app:general}

The word argument works over any finite alphabet. We describe the
changes to the tree semantics, satisfiability tests, and reader
compilation. They give the following decision theorem.

\begin{theorem}\label{thm:general}
For one constructor of any arity and variance, the $\Sigma_1$ theory
over $\BFT$ is decidable over arbitrary, regular, and finite trees.
For $\BF$ the same holds with a covariant constructor of any arity.
For each of these signatures, the arbitrary and regular theories agree.
The compiled small-hole bound in \cref{thm:short} holds for every
arity and variance over $\BFT$.
\end{theorem}

\subsection{Arity and variance}

For arity $n\ge0$, paths use the $n$ argument positions. With variance,
compare labels in reverse order at a common position if its path
has passed through an odd number of contravariant arguments.
Call this parity the polarity of the path. Flipping $\Bot$ and
$\Top$ at odd paths is an involution $\eta$ on trees that sends
the subtype order to the covariant order. This follows at every
common position, since the flip reverses the label order. The map
preserves finiteness and regularity. A subtree in the image is
the image of an original subtree or its leaf-label dual,
according to the polarity of the path to it.

Use the leaf cones of $\eta(t)$. At a child position $i$, cone
membership exchanges signs exactly when argument $i$ is
contravariant. Thus the constructor rows in \cref{sec:readers}
only change the sign of the child variable on contravariant
arguments. The moves still satisfy sign reversal, and a pop
still starts at a signed variable whose root is forbidden.
The proofs of root saturation, realization, and compilation
therefore apply. The same two components over signed variables
give the same action and depth counts. The proof for regular trees
uses residuals of the two cones, exchanging them at odd positions.
Flattening and term evaluation work with any arity, including
zero, as all tree classes are closed under the constructor.

For satisfiability with variance, use a covariant constraint
on two signed copies of each variable. In the first copy of a
literal take the positive variable for every root. In a
constructor literal use the negative child variable on
contravariant arguments and the positive child variable
elsewhere. Form the second copy by exchanging signs, reversing
the inequalities, and exchanging $\Bot$ and $\Top$ in constant
equalities. For example, a lower constructor literal gives
\[
  f(x_1^{s_1},\dots,x_n^{s_n})\le y^+,\qquad
  y^-\le f(x_1^{\bar s_1},\dots,x_n^{\bar s_n}),
\]
where $s_i=-$ on contravariant arguments. Write $\Gamma^\pm$
for the translated constraint. Given a solution $\rho$,
assign $\eta(\rho(x))$ to $x^+$ and its dual to $x^-$.
The ground dual exchanges all leaf labels and reverses the
covariant order. The signed literals hold, because at child $i$
normalization takes the image under $\eta$ and further applies
the dual exactly for a contravariant argument. Conversely,
a solution of $\Gamma^\pm$ with the two signed values dual to
each other decodes by applying $\eta$ to the positive values.

To construct such a solution, let $B$ be the least-shape solution
of the covariant system from \cref{app:shape}. Let $B^\#$ exchange
signed variables and dualize their trees. The translation is
closed under the corresponding operation on literals, so
$B^\#$ also solves it. The two trees $B(p)$ and $B^\#(p)$
have the same constructor positions. Indeed, flipping the
signs in a derivation of a variable inequality and reversing
its endpoints preserves the rules, including decomposition.
Induction on paths then exchanges the lower and upper bounds.
It exchanges the two constructor supports in the label rule
for the least-shape assignment, so the constructor labels
and the domains of the two trees agree.

At every path of a signed variable, select the label from
$B$ when its sign bit plus the path polarity is even, and
from $B^\#$ when it is odd. Here the positive sign has bit
zero. The common shape makes this a tree, and the opposite
sign uses the opposite choice, giving dual trees. The
selection preserves every literal. A variable inequality
uses the same sign on both sides and hence the same choice
at each path. In a constructor literal the choice at a path
$iw$ of the root agrees with the choice at $w$ of the signed
child, since their sign bits differ by the variance at $i$.
Comparison at the root uses the fixed label $f$ on the
constructor side. Constant equalities hold because both
assignments already satisfy them. This gives a solution
with dual signed values, and the domain is unchanged from
$B$. Consequently the label and cycle tests of
\cref{app:classes} decide arbitrary and finite
satisfiability after translation.

For finite transfer with variance, the observations are still bounds
by finite trees, by applying the covariant observation construction
to $\eta(t)$ and then undoing $\eta$. In the signed system impose
each bound on the positive variable after normalization and impose
its dual on the negative variable. For example $g\le x$ imposes
$\eta(g)\le x^+$ and $x^-\le\eta(g)^*$. The signed system has a
finite solution with dual values by assumption, and it has an
arbitrary solution meeting all signed bounds. The covariant
pinning proof gives a finite solution $A$ meeting the signed
bounds. Its sign dual $A^\#$ also meets them. Let $I$ be
any finite solution with dual values supplied by the original
finite solution.

Take the componentwise median of $A,A^\#,I$. To define the
median of three covariant trees, put a word in a given cone
when it belongs to that cone in at least two trees. These
cones are extension closed and disjoint, and thus define a
tree. The median is monotone by the cone inclusions and
commutes with constructor application, because child cones
are selected independently. It commutes with duality and
returns $t$ when two inputs equal $t$. If the trees are
finite, any path beyond their maximum depth has reached a
leaf in each tree and hence a majority cone, so the result
is finite. It follows that the componentwise median of
the three solutions solves the signed literals. Any ground
bound satisfied by two inputs is preserved by monotonicity
and the repeated-input identity. The median is fixed by
sign duality, which interchanges $A$ and $A^\#$ and fixes
$I$. Decoding gives a finite solution of the original
constraint with all the observations. This proves the
finite transfer for arbitrary variance.

\subsection{Removing top}

With no $\Top$ and with covariance, a tree is determined by its
prefix-closed set of constructor positions. The empty set gives
$\Bot$. Inclusion rules out a comparison of a constructor
with a $\Bot$-leaf at the same position. If a constructor position
in the first set is absent from the second, its shortest absent
prefix is a common position labeled $f$ on the left and $\Bot$ on
the right. Thus the order is inclusion. The same set
embeds in the full signature as the very same tree. This embedding
preserves and reflects the order and the constructor. A tree
then has finitely many positions exactly when its constructor
set is finite, because each constructor has finitely many
children. It is regular exactly when the set has finitely many
residuals by paths. Residuals at constructor positions describe
the subtrees there, and all other residuals are empty.

To exclude $\Top$ inside the full signature, add a fresh
variable $g$ and the literals
\[
         g\le f(g,\dots,g),\qquad x\le g\quad(x\in V).
\]
Any positive cone word of the value of $g$ must be nonempty,
because $g$ is below a constructor. The first inequality sends
each such word, with its first letter removed, into the same
positive cone. Iterating would put the empty word there.
Thus $g$ has empty positive cone, and so does every value below
it. Every solution of the guarded constraint therefore restricts
to a solution without top, preserving the query values.

Conversely, for arbitrary or regular trees use the complete
tree with $f$ at every path as the guard. It is regular and
contains every constructor position. In the finite class take
a complete finite tree with $\Bot$ at its frontier, deep enough
to contain the constructor positions of every variable. Its
constructor set is included in that of one more complete level,
so it satisfies the guard inequality. These choices also apply
at arity zero, where the complete constructor tree is $f()$.
This proves the equivalence for entailment in all three classes.
Flattening in the restricted language uses just its constants
and constructor. The guarded reduction, the full-signature
deciders, and Boolean evaluation now prove \cref{thm:general}.
